%!TEX program = uplatex
\RequirePackage{plautopatch}

\documentclass[uplatex,dvipdfmx]{jlreq}

% 余白調整（1ページ収め）
\usepackage[top=18mm,bottom=20mm,left=20mm,right=20mm]{geometry}

\usepackage{amsmath,amssymb}
\usepackage{url}

\pagestyle{empty}

% 行間少し詰める
\renewcommand{\baselinestretch}{0.96}

%--------------------------------
% 講演マクロ
%--------------------------------

\newcommand{\talk}[3]{%
\noindent
#1\quad #2\par
\hspace*{2em}#3\par\smallskip
}

%--------------------------------
% タイトル情報
%--------------------------------

\newcommand{\EwMnumber}{第33回}
\newcommand{\EwMtitle}{双曲力学系\\[2mm]\large 安定性と混沌}
\newcommand{\EwMdate}{2005年2月4日（金）14:30～18:00 ～ 2月5日（土）10:30～17:00}
\newcommand{\EwMplace}{東京都文京区春日1-13-27 中央大学理工学部}

%--------------------------------
% タイトル表示
%--------------------------------

\newcommand{\EwMtitleblock}{

\vspace*{-5mm}

\begin{center}

{\Large ENCOUNTER with MATHEMATICS}

\vspace{2mm}

{\large 数学との遭遇}

\vspace{4mm}

{\Large \bfseries \EwMnumber\ \EwMtitle}

\vspace{4mm}

\EwMdate

\vspace{2mm}

於：\EwMplace

\end{center}

\vspace{2mm}

}

\begin{document}

\EwMtitleblock

%--------------------------------
% プログラム
%--------------------------------

\section*{プログラム}

\subsection*{2月4日（金）}

\talk
{14:30～16:00}
{双曲力学系}
{国府 寛司 氏（京大・理）}

\talk
{16:30～18:00}
{力学系の安定性と通有性 I}
{林 修平 氏（東大・数理）}

\subsection*{2月5日（土）}

\talk
{10:30～11:20}
{ホモクリニック分岐 I}
{浅岡 正幸 氏（京大・理）}

\talk
{11:40～12:50}
{エノン写像}
{三波 篤郎 氏（北見工大）}

\talk
{14:30～15:20}
{ホモクリニック分岐 II}
{浅岡 正幸 氏（京大・理）}

\talk
{15:40～17:00}
{力学系の安定性と通有性 II}
{林 修平 氏（東大・数理）}

\bigskip

17:15～　懇親会（ワイン・パーティー）

%--------------------------------
% 案内文
%--------------------------------

\bigskip

別紙の趣旨に沿った集会の第33回を以上のような予定で開催いたします。  
非専門家向けに入門的な講演をお願い致しました。  
多くの方々のご参加をお待ちしております。  
講演者による講演内容へのご案内を別紙に添付いたしますのでご覧下さい。

%--------------------------------
% 連絡先
%--------------------------------

\section*{連絡先}

112-8551 東京都文京区春日1-13-27 中央大学理工学部数学教室 

tel : 03-3817-1745  

ENCOUNTER with MATHEMATICS  
e-mail : encmath@math.chuo-u.ac.jp  
homepage : \url{http://www.math.chuo-u.ac.jp/ENCwMATH}  

三松 佳彦 : yoshi@math.chuo-u.ac.jp / 高倉 樹 : takakura@math.chuo-u.ac.jp

\end{document}
